\section{Estado del Arte}

\subsubsection{Antecedentes Académicos y Técnicos}

La literatura relacionada con sistemas digitales para restaurantes, omnicanalidad, pedidos
en mesa, men\'us QR y pagos digitales muestra que la digitalizaci\'on puede generar mejoras
operativas importantes, aunque tambi\'en introduce retos asociados con la experiencia del
usuario, la integraci\'on de canales y la complejidad t\'ecnica del sistema. Los antecedentes
se organizan en cinco ejes:

\noindent\textbf{Digitalizaci\'on, omnicanalidad y desempe\~no operativo.}
Diversos trabajos recientes han analizado el efecto de la transformaci\'on digital en la
industria restaurantera. Lee et~al. (2024) documentan que la digitalizaci\'on del
\textit{front of house} y del \textit{back of house} contribuye a fortalecer el desempe\~no
organizacional, especialmente cuando los sistemas permiten integrar informaci\'on de distintas
\'areas del negocio \cite{katagri_smart_2025}. De forma complementaria, Wang et~al. (2024) proponen modelos
anal\'iticos para estudiar la operaci\'on restaurantera en contextos omnicanal mediante teor\'ia
de colas y teor\'ia de juegos, mostrando que la conveniencia de operar mediante canales
presenciales, digitales o una combinaci\'on de ambos depende de factores como la sensibilidad
al tiempo de espera y la capacidad del restaurante \cite{oracle_consumer_study}.

Katagri y Balavalad (2025) presentaron un sistema inteligente de ordenamiento digital para
restaurantes que integra una interfaz responsiva, actualizaciones en tiempo real y pagos
seguros, reportando reducciones en tiempos de error y procesamiento de \'ordenes al eliminar
la intermediaci\'on manual \cite{katagri_smart_2025}. Sus resultados respaldan la pertinencia de dise\~nar una
plataforma que integre la operaci\'on presencial del restaurante con mecanismos digitales
de interacci\'on, manteniendo una base de datos centralizada.

\noindent\textbf{Autoservicio en mesa y aplicaciones m\'oviles.}
Tan y Netessine (2020) reportan que las tecnolog\'ias de autoservicio en mesa pueden incrementar
las ventas por cuenta, reducir la duraci\'on del servicio y mejorar la productividad al aumentar
la rotaci\'on de mesas \cite{lee_digital_2024}. No obstante, Xu et~al. (2024) se\~nalan que la incorporaci\'on
intensiva de aplicaciones m\'oviles de pedido tambi\'en puede relacionarse con disminuciones en
ciertos indicadores de satisfacci\'on del consumidor \cite{xu_mobile_2024}, lo que sugiere que la mejora operativa
no siempre se traduce autom\'aticamente en una mejor experiencia del usuario. Estos hallazgos
orientan el dise\~no del presente proyecto hacia interfaces de m\'inima fricci\'on y flujos guiados.

\noindent\textbf{C\'odigos QR en restaurantes.}
Koay y Ang (2024) examinan la intenci\'on de uso de men\'us QR a partir de variables como
utilidad percibida, facilidad de uso, influencia social y percepci\'on de privacidad,
encontrando que la utilidad percibida y los h\'abitos previos son los predictores m\'as fuertes
de adopci\'on \cite{koay_qr_2024}. Sin embargo, Hewage et~al. (2024) muestran que los men\'us QR pueden
generar inconveniencia percibida y afectar negativamente la lealtad del cliente cuando existe
alta preferencia por la interacci\'on con el personal de servicio \cite{hewage_qr_2024}. Yigito\u{g}lu et~al. (2025)
destacan que el riesgo percibido y el contexto cultural modifican la aceptaci\'on de estas
tecnolog\'ias \cite{yigitoglu_sustainable_2025}. Estos hallazgos son especialmente relevantes para este trabajo, ya que
el dise\~no del sistema deber\'a priorizar facilidad de uso, rapidez de acceso y baja fricci\'on
durante la interacci\'on del comensal con el c\'odigo QR de la mesa virtual.

\noindent\textbf{Sincronizaci\'on en tiempo real y arquitectura.}
Iovescu y Tudose (2024) proponen arquitecturas orientadas a eventos para sistemas que requieren
actualizaci\'on casi inmediata entre m\'ultiples actores, destacando la importancia de la
idempotencia en comandos y transacciones \cite{iovescu_realtime_2024}. Sus aportaciones son directamente aplicables
al presente sistema, donde la confirmaci\'on de una orden o de un pago no debe provocar
duplicidad de registros si la misma operaci\'on se procesa m\'as de una vez por error de red o
reconexi\'on.

El sistema de pedidos QR desarrollado por Choo et~al. (2023) publicado en el
\textit{Journal of Informatics and Web Engineering} integra generaci\'on de QR por mesa,
interfaz de administraci\'on de men\'u, facturaci\'on automatizada y visualizaci\'on de datos
anal\'iticos, reportando una reducci\'on significativa en errores de captura respecto a sistemas
con comandas en papel \cite{choo_qr_2023}. Si bien este trabajo no aborda la divisi\'on de cuenta entre
comensales ni la arquitectura serverless, sienta una base s\'olida para la presente propuesta.

\noindent\textbf{Pagos digitales, propinas y divisi\'on de cuenta.}
Alsaati et~al. (2026) muestran que la percepci\'on de manipulaci\'on en las interfaces de
pago puede afectar negativamente la satisfacci\'on del usuario y la favorabilidad hacia la
marca, lo que resalta la importancia de dise\~nar flujos de pago transparentes \cite{pci_dss_4}.
Desde el punto de vista algor\'itmico, Bultel et~al. (2020) demuestran que algunos problemas
de minimizaci\'on de transacciones en contextos de reparto de gastos pueden ser
computacionalmente complejos \cite{bultel_faster_2020}, lo cual resulta relevante como antecedente te\'orico
para el m\'odulo de divisi\'on de cuentas del presente proyecto.

\noindent\textbf{Anal\'itica de datos y procesamiento masivo.}
El procesamiento de grandes vol\'umenes de datos en la industria restaurantera permite
identificar patrones de consumo, optimizar inventarios y predecir demanda. Zaharia et~al. (2012)
introducen el concepto de conjuntos de datos distribuidos resilientes (RDD), permitiendo que
motores como Apache Spark realicen an\'alisis en memoria con alta tolerancia a fallos y
escalabilidad horizontal \cite{zaharia2012resilient}. Kleppmann (2017) destaca que la transici\'on de
sistemas puramente transaccionales (OLTP) hacia arquitecturas h\'ibridas con capas
anal\'iticas (OLAP) es fundamental para la escalabilidad de aplicaciones intensivas en
datos \cite{kleppmann2017designing}. En el presente proyecto, esta base te\'orica sustenta el uso de Spark
sobre una arquitectura Medallion para procesar las transacciones operativas y generar
tableros de inteligencia de negocio estrat\'egica.

La Tabla~\ref{tab:antecedentes} resume los antecedentes acad\'emicos y t\'ecnicos
m\'as relevantes para el presente trabajo.

{\small
\setlength{\tabcolsep}{3pt}
\renewcommand{\arraystretch}{1.5}
\begin{longtable}{>{\raggedright\arraybackslash}p{3.5cm}
                >{\centering\arraybackslash}p{1.2cm}
                >{\raggedright\arraybackslash}p{4.5cm}
                >{\raggedright\arraybackslash}p{6.5cm}}
\caption{Resumen de Antecedentes Académicos y Técnicos Relacionados con el Proyecto} \label{tab:antecedentes} \\
\toprule
\thead{Trabajo} & \thead{A\~no} & \thead{Aporte principal} & \thead{Aplicabilidad al proyecto} \\
\midrule
\endfirsthead
\caption[]{TABLA \thetable{} (Continuación)} \\
\toprule
\thead{Trabajo} & \thead{A\~no} & \thead{Aporte principal} & \thead{Aplicabilidad al proyecto} \\
\midrule
\endhead
\bottomrule
\endfoot
Lee, Jang y Kim \cite{lee_digital_2024} & 2024 &
  Relaci\'on positiva entre transformaci\'on digital y productividad en restaurantes. &
  Justifica la integraci\'on tecnol\'ogica de procesos operativos. \\
\rowcolor{altrow}
Wang et al. \cite{wang_omnichannel_2024} & 2024 &
  Modela la conveniencia de operar en canales dine-in, delivery y omnicanal. &
  Sustenta el enfoque omnicanal del sistema. \\
Katagri y Balavalad \cite{katagri_smart_2025} & 2025 &
  Sistema inteligente de pedidos con interfaz responsiva y pagos seguros. &
  Referente directo para la arquitectura de pedidos y reducci\'on de errores. \\
\rowcolor{altrow}
Tan y Netessine \cite{tan_tabletop_2020} & 2020 &
  Tecnolog\'ias en mesa aumentan ventas y reducen duraci\'on del servicio. &
  Apoya el uso de pedidos y pagos digitales en mesa. \\
Koay y Ang \cite{koay_qr_2024} & 2024 &
  Analiza intenci\'on de uso de men\'us QR con variables UTAUT2 y privacidad. &
  Orienta el dise\~no UX del acceso por QR. \\
\rowcolor{altrow}
Hewage, Boman y Lefebvre \cite{hewage_qr_2024} & 2024 &
  Men\'us QR pueden reducir lealtad por inconveniencia percibida. &
  Advierte sobre la necesidad de minimizar fricci\'on y ofrecer alternativas. \\
Yigito\u{g}lu et al. \cite{yigitoglu_sustainable_2025} & 2025 &
  Impacto de men\'us QR sustentables en calidad de servicio y satisfacci\'on. &
  Fundamenta el dise\~no accesible y de bajo riesgo percibido. \\
\rowcolor{altrow}
Iovescu y Tudose \cite{iovescu_realtime_2024} & 2024 &
  Arquitecturas orientadas a eventos para sincronizaci\'on en tiempo real. &
  Fundamenta el uso de eventos e idempotencia para \'ordenes y pagos. \\
Choo et al. \cite{choo_qr_2023} & 2023 &
  Sistema QR con men\'u digital, facturaci\'on automatizada y anal\'itica. &
  Antecedente directo del m\'odulo de c\'odigos QR y tablero de administraci\'on. \\
\rowcolor{altrow}
Alsaati et al. \cite{alsaati_etipping_2026} & 2026 &
  La manipulaci\'on percibida en interfaces de propina afecta satisfacci\'on. &
  Sustenta dise\~no \'etico y transparente del flujo de pagos. \\
\rowcolor{altrow}
Zaharia et al. \cite{zaharia2012resilient} & 2012 &
  Arquitectura de procesamiento distribuido en memoria (RDD). &
  Base tecnol\'ogica para el pipeline anal\'itico de Spark. \\
Bultel et al. \cite{bultel_faster_2020} & 2020 &
  Minimizar transacciones en reparto de gastos puede ser NP-completo. &
  Antecedente te\'orico para modelos de divisi\'on de pagos. \\
\end{longtable}
}

\subsubsection{Soluciones de Mercado y Aplicaciones Comerciales}

En el mercado actual existen diversas plataformas orientadas a la gesti\'on de restaurantes.
Para el an\'alisis comparativo del presente trabajo, las soluciones se organizaron en tres grupos:
plataformas POS y KDS integradas; plataformas de gesti\'on omnicanal de \'ordenes; y pasarelas
o procesadores de pago con capacidades de integraci\'on digital.

\noindent\textbf{Plataformas POS y KDS integradas.}
Soluciones como Toast, Square~for~Restaurants, Lightspeed~Restaurant y TouchBistro incorporan
funcionalidades de \textit{split bill}, pagos m\'ultiples, visualizaci\'on de \'ordenes en cocina y
herramientas para mejorar la eficiencia operativa \cite{toast_contactless_2026}--\cite{touchbistro_pos}. En general, estas plataformas
muestran un alto nivel de madurez en la gesti\'on interna del restaurante, particularmente en
procesos administrados por meseros, caja o cocina. No obstante, en varios casos la experiencia
de pago dividido sigue estando m\'as centrada en la l\'ogica operativa del establecimiento que en
un flujo completamente aut\'onomo para el comensal, y muchas forman parte de ecosistemas
cerrados con dependencia de hardware o configuraciones espec\'ificas del proveedor.

CoverAtTable (CoverManager) ofrece un sistema con c\'odigo QR en mesa para pedidos y pagos
que permite al comensal pagar la cuenta o realizar el pedido desde su dispositivo m\'ovil sin
llamar al camarero \cite{mdn_websocket}. Sin embargo, no integra la gesti\'on de mesas virtuales compartidas con
divisi\'on aut\'onoma de cuenta entre m\'ultiples comensales identificados individualmente.

\noindent\textbf{Plataformas omnicanal de gesti\'on de \'ordenes.}
Soluciones como Deliverect destacan por su capacidad de integraci\'on con POS existentes y
por coordinar men\'us, estados y flujos de \'ordenes entre distintos or\'igenes. Sin embargo,
funcionan m\'as como orquestadores de canales que como una soluci\'on integral de interacci\'on
comensal-mesa-pago.

\noindent\textbf{Pasarelas y procesadores de pago.}
Stripe y Conekta ofrecen infraestructura cr\'itica para el procesamiento de cobros digitales,
tokenizaci\'on y conciliaci\'on de transacciones \cite{pci_dss_4, stripe_pricing}. Estas soluciones se consideran
complementarias al sistema propuesto: no resuelven por s\'i mismas la l\'ogica de divisi\'on de
cuenta en una mesa virtual, pero proporcionan la base tecnol\'ogica necesaria para ejecutar
los pagos de forma segura y trazable.

La Tabla~\ref{tab:soluciones_comerciales} presenta una comparaci\'on de las soluciones
comerciales m\'as relevantes identificadas en el estado del arte.

\begin{longtable}{>{\raggedright\arraybackslash}p{2.2cm}
                >{\raggedright\arraybackslash}p{2.5cm}
                >{\centering\arraybackslash}p{1.3cm}
                >{\centering\arraybackslash}p{1.3cm}
                >{\centering\arraybackslash}p{1.3cm}
                >{\raggedright\arraybackslash}p{4.1cm}}
\caption{Comparación de Soluciones Comerciales Relacionadas con el Proyecto}
\label{tab:soluciones_comerciales}\\
\toprule
\thead{Proveedor} & \thead{Producto} & \thead{QR / mesa} &
\thead{Split bill} & \thead{Mesa virtual} & \thead{Observaciones} \\
\midrule
\endfirsthead

\toprule
\thead{Proveedor} & \thead{Producto} & \thead{QR / mesa} &
\thead{Split bill} & \thead{Mesa virtual} & \thead{Observaciones} \\
\midrule
\endhead

\midrule
\multicolumn{6}{r}{Contin\'ua en la siguiente p\'agina} \\
\endfoot

\bottomrule
\endlastfoot
Toast & POS + Mobile Order \& Pay + KDS & S\'i & S\'i & No &
  Ecosistema robusto, dependencia del stack del proveedor. \\
\rowcolor{altrow}
Square & Square for Restaurants + KDS & S\'i & S\'i & No &
  Soporta pagos m\'ultiples y divisi\'on por \'items con integraci\'on propia. \\
Lightspeed & Restaurant + Order Anywhere + KDS & S\'i & S\'i & No &
  Buen soporte para cuenta dividida y pedidos por QR, con variaciones por regi\'on. \\
\rowcolor{altrow}
CoverManager & CoverAtTable & S\'i & Parcial & No &
  QR en mesa para pedidos y pagos, sin gesti\'on de comensales individuales. \\
Deliverect & Order Management + Integraciones & Parcial & No central & No &
  Destaca en agregaci\'on omnicanal, no en mesa y pago dividido. \\
\rowcolor{altrow}
Mercado Pago & Checkout / Links / QR & No espec\'ifico & Parcial & No &
  \'Util como capa de pago, no como gestor integral de restaurante. \\
Stripe & Payments + Connect & No espec\'ifico & Parcial & No &
  Flexible por API; la l\'ogica de divisi\'on debe implementarse en la app. \\
\rowcolor{altrow}
Biomenus & Men\'u digital PWA con nutrici\'on, alerg\'enos y filtros & S\'i & No & No &
  Especializado en informaci\'on nutricional y filtros diet\'eticos; sin KDS, sin cuentas por comensal ni gesti\'on de mesas. \\
Loyverse POS & POS + KDS + inventario (freemium) & Parcial* & Parcial & No &
  KDS y divisi\'on de tickets gratuitos; sin mapeo de mesas, sin QR de sesiones, sin cuentas aut\'onomas por comensal. QR ordering s\'olo v\'ia integraciones de terceros. \\
\rowcolor{altrow}
\textbf{El presente sistema} & PWA omnicanal & \textbf{S\'i} & \textbf{S\'i} &
  \textbf{S\'i} & Mesa virtual con contrase\~na, divisi\'on aut\'onoma de cuenta y arquitectura serverless. \\
\end{longtable}

\noindent\footnotesize{*QR ordering en Loyverse requiere integraciones externas del App Marketplace (ej. MENU~TIGER, Fuudey).}

\normalsize

\subsubsection{Análisis Detallado: Biomenus}

Biomenus es una plataforma de men\'us digitales inteligentes especializada en la presentaci\'on
de informaci\'on nutricional, declaraci\'on de al\'ergenos y filtros diet\'eticos (vegano,
sin gluten, sin l\'acteos, entre otros) \cite{biomenus}. Su propuesta de valor se centra en empoderar al
comensal con informaci\'on detallada sobre los platillos antes de ordenar, facilitando la
toma de decisiones para personas con restricciones alimentarias o necesidades espec\'ificas
de salud.

La plataforma permite a los restaurantes crear y publicar men\'us digitales accesibles
mediante c\'odigo QR, con soporte para actualizaci\'on en tiempo real de precios, fotos,
descripciones y disponibilidad de platillos. Destaca por ofrecer perfiles nutricionales
completos por plato y la posibilidad de que los comensales filtren el cat\'alogo seg\'un sus
preferencias diet\'eticas.

Sin embargo, Biomenus opera exclusivamente como una herramienta de \textit{consulta
de men\'u}, sin incluir funcionalidades de gesti\'on de pedidos, seguimiento de \'ordenes
en cocina (KDS), gesti\'on de mesas, ni divisi\'on de cuentas entre comensales. El flujo
de pago queda completamente fuera de su alcance. En este sentido, constituye una soluci\'on
complementaria pero no sustituta de un sistema integral de gesti\'on restaurantera.

La Tabla~\ref{tab:comp_biomenus} presenta el an\'alisis comparativo entre Biomenus y la
propuesta del presente proyecto.

{\small
\setlength{\tabcolsep}{3pt}
\renewcommand{\arraystretch}{1.5}
\begin{longtable}{>{\raggedright\arraybackslash}p{4.5cm}
                >{\centering\arraybackslash}p{2.5cm}
                >{\centering\arraybackslash}p{2.5cm}
                >{\raggedright\arraybackslash}p{6cm}}
\caption{Análisis Comparativo: Biomenus vs. Sistema Propuesto} \label{tab:comp_biomenus} \\
\toprule
\thead{Capacidad} & \thead{Biomenus} & \thead{Sistema} & \thead{Observaci\'on} \\
\midrule
\endfirsthead
\caption[]{TABLA \thetable{} (Continuación)} \\
\toprule
\thead{Capacidad} & \thead{Biomenus} & \thead{Sistema} & \thead{Observaci\'on} \\
\midrule
\endhead
\bottomrule
\endfoot
Men\'u digital con QR & S\'i & S\'i & Ambas ofrecen acceso al men\'u v\'ia QR sin app nativa. \\
\rowcolor{altrow}
Informaci\'on nutricional y al\'ergenos & \textbf{S\'i (detallada)} & No (v1.0) & Biomenus supera al sistema propuesto en profundidad de informaci\'on diet\'etica; identificado como trabajo futuro. \\
Filtros diet\'eticos (vegano, gluten-free) & \textbf{S\'i} & No (v1.0) & Funcionalidad de roadmap en el sistema. \\
\rowcolor{altrow}
Actualizaci\'on de men\'u en tiempo real & S\'i & S\'i & Ambas permiten cambios inmediatos sin reimprimir. \\
Gesti\'on de \'ordenes (comensal ordena) & No & \textbf{S\'i} & El sistema incluye flujo completo de canasta y confirmaci\'on. \\
\rowcolor{altrow}
KDS para Chef en tiempo real & No & \textbf{S\'i} & El sistema sincroniza \'ordenes con el tablero del Chef v\'ia Supabase Realtime. \\
Mesa virtual con contrase\~na & No & \textbf{S\'i} & Diferenciador central del sistema: acceso seguro por JWT + contrase\~na. \\
\rowcolor{altrow}
Divisi\'on de cuenta por comensal & No & \textbf{S\'i} & El sistema permite liquidar cuentas individuales y mesa compartida de forma aut\'onoma. \\
Dashboard administrativo & No & \textbf{S\'i} & El sistema incluye m\'etricas por cadena, zona y sucursal. \\
\rowcolor{altrow}
Gesti\'on de roles y personal & No & \textbf{S\'i} & El sistema implementa jerarqu\'ia de 5 roles con RLS en Supabase. \\
Modelo de negocio & SaaS / freemium & Desarrollo propio & Biomenus opera como servicio externo; el sistema es propiedad total del c\'odigo. \\
\end{longtable}
}

\subsubsection{Análisis Detallado: Loyverse POS}

Loyverse~POS es un sistema de punto de venta (POS) en la nube con modelo freemium,
disponible para iOS y Android, orientado principalmente a caf\'es, restaurantes de
servicio r\'apido y comercios minoristas peque\~nos \cite{loyverse_features}. Su plan gratuito incluye funciones
habitualmente de pago en otros sistemas: POS b\'asico, KDS, pantalla de cliente,
programa de lealtad, gesti\'on de inventario y reportes anal\'iticos en tiempo real \cite{fitsmallbusiness_loyverse}.
Las funciones avanzadas (gesti\'on de empleados, inventario avanzado e integraciones) se
ofrecen como m\'odulos adicionales con tarifa mensual por tienda.

En cuanto a la experiencia del comensal, Loyverse no cuenta con QR de sesi\'on nativo;
el ordenamiento por QR s\'olo es posible mediante integraciones del App~Marketplace como
MENU~TIGER o Fuudey \cite{posusa_loyverse}. Estas integraciones registran los pedidos en el sistema
de Loyverse, pero no implementan el concepto de mesa virtual protegida por contrase\~na
ni la gesti\'on aut\'onoma de cuentas individuales por comensal.

Una limitaci\'on documentada de Loyverse es la ausencia de mapeo visual de mesas
(\textit{table mapping}), lo que lo hace poco adecuado para restaurantes de servicio
completo donde la gesti\'on espacial de mesas es esencial \cite{fitsmallbusiness_loyverse, posusa_loyverse}. La divisi\'on de
tickets existe pero opera desde el lado del cajero/mesero, no desde el dispositivo
del comensal.

La Tabla~\ref{tab:comp_loyverse} presenta el an\'alisis comparativo entre Loyverse y la
propuesta del presente proyecto.

{\small
\setlength{\tabcolsep}{3pt}
\renewcommand{\arraystretch}{1.5}
\begin{longtable}{>{\raggedright\arraybackslash}p{4.5cm}
                >{\centering\arraybackslash}p{2.5cm}
                >{\centering\arraybackslash}p{2.5cm}
                >{\raggedright\arraybackslash}p{6cm}}
\caption{Análisis Comparativo: Loyverse vs. Sistema Propuesto} \label{tab:comp_loyverse} \\
\toprule
\thead{Capacidad} & \thead{Loyverse} & \thead{Sistema} & \thead{Observaci\'on} \\
\midrule
\endfirsthead
\caption[]{TABLA \thetable{} (Continuación)} \\
\toprule
\thead{Capacidad} & \thead{Loyverse} & \thead{Sistema} & \thead{Observaci\'on} \\
\midrule
\endhead
\bottomrule
\endfoot
POS (toma de \'ordenes por personal) & \textbf{S\'i (gratuito)} & Limitado & Loyverse supera al sistema en POS tradicional; el sistema se enfoca en autoservicio del comensal. \\
\rowcolor{altrow}
KDS para cocina en tiempo real & \textbf{S\'i (gratuito)} & S\'i & Ambas ofrecen KDS; Loyverse usa app nativa, el sistema usa WebSockets + Supabase Realtime. \\
Programa de lealtad & \textbf{S\'i (gratuito)} & No (v1.0) & Loyverse incluye puntos de fidelizaci\'on; identificado como trabajo futuro en el sistema. \\
\rowcolor{altrow}
Gesti\'on de inventario & \textbf{S\'i (avanzado, pago)} & No & Loyverse gestiona stock e ingredientes; fuera del alcance del sistema v1.0. \\
Reportes de ventas & \textbf{S\'i} & S\'i & Loyverse reporta por empleado, turno e \'item; el sistema por cadena, zona y sucursal. \\
\rowcolor{altrow}
Mapeo visual de mesas & No & No (v1.0) & Ambas carecen de plano de planta en sus versiones actuales. \\
QR nativo para ordenar desde mesa & No nativo* & \textbf{S\'i (JWT)} & El sistema implementa QR de sesi\'on \'unica con contrase\~na; Loyverse requiere app de terceros. \\
\rowcolor{altrow}
Mesa virtual protegida por contrase\~na & No & \textbf{S\'i} & Diferenciador central del sistema: sesiones privadas por grupo de comensales. \\
Divisi\'on aut\'onoma de cuenta (comensal) & No & \textbf{S\'i} & En Loyverse la divisi\'on la hace el cajero; en el sistema la hace el propio comensal. \\
\rowcolor{altrow}
Gesti\'on multi-nivel (cadena/zona/sucursal) & Parcial (multi-tienda) & \textbf{S\'i} & El sistema implementa jerarqu\'ia de 3 niveles con acceso por rol; Loyverse gestiona m\'ultiples tiendas pero sin jerarqu\'ia org. \\
Funcionamiento offline & \textbf{S\'i} & No & Loyverse sincroniza ventas offline; el sistema requiere conexi\'on permanente. \\
\rowcolor{altrow}
Modelo de acceso & App nativa iOS/Android & PWA (navegador) & Loyverse requiere instalaci\'on de app; el sistema es accesible sin instalar nada. \\
Costo base & \textbf{Gratuito} & Infraestructura (\$914~MXN/mes) & Loyverse no tiene costo de software base; el sistema tiene costo de infraestructura cloud. \\
\end{longtable}
}

\noindent\footnotesize{*QR ordering en Loyverse requiere integraciones del App Marketplace (MENU~TIGER, Fuudey, Orderlina).}

\normalsize

\subsubsection{Síntesis del Estado del Arte}

A partir de la revisi\'on acad\'emica y comercial realizada, se observa que existen avances
importantes en digitalizaci\'on de restaurantes, autoservicio en mesa, men\'us QR, plataformas
omnicanal y procesamiento de pagos digitales. Sin embargo, tambi\'en se identifican vac\'ios
relevantes en la integraci\'on de estas capacidades dentro de una soluci\'on unificada centrada
en el comensal.

Los an\'alisis comparativos de Biomenus y Loyverse~POS permiten delimitar con precisi\'on las
ventajas y las brechas del presente proyecto. \textbf{Biomenus} lidera en informaci\'on
nutricional y filtros diet\'eticos, funcionalidades que el sistema no incluye en v1.0 pero que
representan una oportunidad clara de diferenciaci\'on futura. En contraste, Biomenus carece
completamente de gesti\'on de \'ordenes, KDS, cuentas por comensal y dashboard administrativo,
\'areas en las que el sistema ofrece cobertura completa. \textbf{Loyverse~POS} supera al sistema propuesto
en su modelo POS tradicional (gratuito, nativo, con soporte offline y programa de lealtad),
pero no implementa el concepto de mesa virtual protegida por contrase\~na, ni la divisi\'on
aut\'onoma de cuenta desde el dispositivo del comensal \cite{fitsmallbusiness_loyverse}. El QR ordering de
Loyverse depende de integraciones externas y no ofrece la trazabilidad individual por comensal
que distingue a la propuesta del presente proyecto.

En particular, las plataformas comerciales existentes tienden a resolver adecuadamente la
divisi\'on de cuenta desde la perspectiva operativa del restaurante (cajero o mesero como
actor), pero ofrecen menor flexibilidad cuando se busca que los propios comensales gestionen,
en tiempo real, pedidos, divisiones por \'item y contribuciones colaborativas dentro de una
misma mesa virtual con identidad propia.

La literatura muestra que el uso de tecnolog\'ias como men\'us QR y autoservicio puede mejorar
la eficiencia operativa, aunque no garantiza por s\'i mismo una mejor experiencia de usuario
\cite{hewage_qr_2024, nra_state_2023}. Por ello, el dise\~no de la presente soluci\'on requiere equilibrar productividad,
facilidad de uso, transparencia en el estado de cuentas y consistencia en tiempo real.

En este sentido, el presente proyecto se posiciona como una propuesta que busca integrar
en una sola plataforma la gesti\'on de mesas, \'ordenes y liquidaci\'on de cuentas divididas,
combinando acceso mediante c\'odigos QR con JWT de un solo uso, sincronizaci\'on en tiempo
real v\'ia Supabase~Realtime y trazabilidad de transacciones. A diferencia de las soluciones
comerciales analizadas, el sistema incorpora un motor de procesamiento distribuido mediante
Apache Spark bajo una arquitectura Medallion, lo que permite transformar los datos operativos
en inteligencia de negocio estrat\'egica para la gesti\'on multinivel de cadenas restauranteras
en M\'exico.

% ============================================================
